\documentclass[12pt]{article}
\usepackage{amsmath, amssymb, amsthm, graphicx, geometry, enumitem}
\geometry{a4paper, margin=1in}

\title{Homework Section 12.3: The Dot Product\\ \large Comprehensive Problem Set and Solutions}
\author{Generated Study Set}
\date{}

\begin{document}

\maketitle

\section*{Instructions}
This problem set contains 30 questions designed to test your mastery of the Dot Product as detailed in Section 12.3. 
\begin{itemize}
    \item \textbf{Part A}: Conceptual Checks (Meaningful vs. Meaningless)
    \item \textbf{Part B}: Basic Computations and Angles
    \item \textbf{Part C}: Orthogonality and Parallelism
    \item \textbf{Part D}: Geometric Applications (Triangles and Shapes)
    \item \textbf{Part E}: Projections and Decompositions
    \item \textbf{Part F}: Physics Applications (Work)
    \item \textbf{Part G}: Advanced Computations and Variables
\end{itemize}

\hrulefill

\section*{Part A: Conceptual Checks}

\textbf{Problem 1.} Determine whether each of the following expressions is meaningful or meaningless. If meaningful, state whether the result is a scalar or a vector. If meaningless, explain why. Assume $\mathbf{u}, \mathbf{v}, \mathbf{w}$ are vectors in $\mathbb{R}^3$.
\begin{enumerate}[label=(\alph*)]
    \item $(\mathbf{u} \cdot \mathbf{v}) + \mathbf{w}$
    \item $|\mathbf{u}|(\mathbf{v} \cdot \mathbf{w})$
    \item $(\mathbf{u} \cdot \mathbf{v}) \cdot \mathbf{w}$
    \item $(\mathbf{u} + \mathbf{v}) \cdot (\mathbf{u} - \mathbf{v})$
\end{enumerate}

\section*{Part B: Basic Computations and Angles}

\textbf{Problem 2.} Find $\mathbf{a} \cdot \mathbf{b}$ for $\mathbf{a} = \langle 4, 1, \frac{1}{4} \rangle$ and $\mathbf{b} = \langle 6, -3, -8 \rangle$.

\textbf{Problem 3.} Find the dot product of $\mathbf{u} = 2\mathbf{i} - 3\mathbf{j} + 5\mathbf{k}$ and $\mathbf{v} = \mathbf{i} + 2\mathbf{j} - \mathbf{k}$.

\textbf{Problem 4.} Find the angle $\theta$ between vectors $\mathbf{a} = \langle 3, 4 \rangle$ and $\mathbf{b} = \langle 5, 12 \rangle$. (Give an exact answer involving arccos).

\textbf{Problem 5.} Find the angle between $\mathbf{a} = \langle 1, -2, 3 \rangle$ and $\mathbf{b} = \langle -3, 6, -9 \rangle$. Explain the geometric significance of your result.

\section*{Part C: Orthogonality and Parallelism}

\textbf{Problem 6.} Determine whether $\mathbf{a} = \langle 6, -2, 4 \rangle$ and $\mathbf{b} = \langle -3, 1, -2 \rangle$ are orthogonal, parallel, or neither.

\textbf{Problem 7.} Determine whether $\mathbf{u} = \langle 5, 3, 7 \rangle$ and $\mathbf{v} = \langle 6, -8, -2 \rangle$ are orthogonal, parallel, or neither.

\textbf{Problem 8.} Find the value(s) of $x$ such that $\mathbf{a} = \langle x, x, 2 \rangle$ and $\mathbf{b} = \langle 1, x, -3 \rangle$ are orthogonal.

\textbf{Problem 9.} Find the value of $b$ such that vectors $\langle -2, b, 4 \rangle$ and $\langle 4, -8, -8 \rangle$ are parallel.

\section*{Part D: Geometric Applications}

\textbf{Problem 10.} Consider the triangle with vertices $A(2, -1, 0)$, $B(4, 1, 1)$, and $C(4, -5, 4)$. Find the angle at vertex $A$ ($\angle BAC$) to the nearest degree.

\textbf{Problem 11.} Use vectors to determine if the triangle with vertices $P(1, -3, -2)$, $Q(2, 0, -4)$, and $R(6, -2, -5)$ is a right-angled triangle. If so, at which vertex is the right angle?

\textbf{Problem 12.} Find the direction cosines and direction angles for the vector $\mathbf{v} = \langle 1, 2, 2 \rangle$.

\textbf{Problem 13.} Two vectors $\mathbf{a}$ and $\mathbf{b}$ form a rhombus. If $\mathbf{a} = \langle 1, 4 \rangle$ and the side lengths are equal, verify algebraically using the dot product that the diagonals are perpendicular. (Hint: The diagonals are $\mathbf{a}+\mathbf{b}$ and $\mathbf{a}-\mathbf{b}$, assume $|\mathbf{a}| = |\mathbf{b}|$).

\textbf{Problem 14.} Three vectors form a cycle (a closed triangle) such that $\mathbf{u} + \mathbf{v} + \mathbf{w} = \mathbf{0}$. If $|\mathbf{u}|=3$, $|\mathbf{v}|=4$, and $|\mathbf{w}|=5$, find $\mathbf{u} \cdot \mathbf{v}$. (Hint: $\mathbf{w} = -(\mathbf{u}+\mathbf{v})$ and square both sides).

\section*{Part E: Projections and Decompositions}

\textbf{Problem 15.} Find the scalar projection of $\mathbf{b} = \langle 4, 2, 0 \rangle$ onto $\mathbf{a} = \langle 1, 1, 1 \rangle$.

\textbf{Problem 16.} Find the vector projection of $\mathbf{b} = \langle -2, 3, 1 \rangle$ onto $\mathbf{a} = \langle 2, 5, -1 \rangle$.

\textbf{Problem 17.} Decompose vector $\mathbf{u} = \langle 3, 4 \rangle$ into two vectors $\mathbf{u}_{\parallel}$ and $\mathbf{u}_{\perp}$, where $\mathbf{u}_{\parallel}$ is parallel to $\mathbf{v} = \langle 1, 1 \rangle$ and $\mathbf{u}_{\perp}$ is orthogonal to $\mathbf{v}$.

\textbf{Problem 18.} If $\text{proj}_{\mathbf{a}} \mathbf{b} = \langle 2, 1 \rangle$ and $\mathbf{a} = \langle 4, 2 \rangle$, is this result consistent? Explain why or why not using the definition of vector projection.

\textbf{Problem 19.} Find the orthogonal component of $\mathbf{b}$ with respect to $\mathbf{a}$ ($\text{orth}_{\mathbf{a}}\mathbf{b}$) given $\mathbf{a} = \langle 2, 0, 1 \rangle$ and $\mathbf{b} = \langle 4, 1, 2 \rangle$.

\section*{Part F: Physics Applications (Work)}

\textbf{Problem 20.} A constant force $\mathbf{F} = \langle 5, 2, 6 \rangle$ (Newtons) moves an object from point $P(1, -1, 2)$ to $Q(3, 4, -1)$ (meters). Calculate the work done.

\textbf{Problem 21.} A toy wagon is pulled a distance of 100 meters along a horizontal path by a constant force of 50 Newtons. The handle of the wagon is held at an angle of $60^\circ$ above the horizontal. Find the work done.

\textbf{Problem 22.} A force $\mathbf{F} = 3\mathbf{i} - \mathbf{j} + \mathbf{k}$ moves a particle along the line segment from origin to point $(2, 2, 2)$. Is the work done positive, negative, or zero? Interpret the result physically.

\section*{Part G: Advanced Computations and Variables}

\textbf{Problem 23.} Find a unit vector that is orthogonal to both $\mathbf{a} = \langle 1, 1, 0 \rangle$ and $\mathbf{b} = \langle 1, -1, 1 \rangle$.

\textbf{Problem 24.} Find the value of $x$ such that the angle between $\mathbf{a} = \langle 1, 0, 1 \rangle$ and $\mathbf{b} = \langle x, 1, 0 \rangle$ is $60^\circ$.

\textbf{Problem 25.} Use the definition of the dot product to prove the Parallelogram Law: $|\mathbf{a} + \mathbf{b}|^2 + |\mathbf{a} - \mathbf{b}|^2 = 2|\mathbf{a}|^2 + 2|\mathbf{b}|^2$.

\textbf{Problem 26.} If $|\mathbf{u}| = 3$, $|\mathbf{v}| = 5$, and the dot product $\mathbf{u} \cdot \mathbf{v} = -10$, find the magnitude of the vector $2\mathbf{u} - \mathbf{v}$.

\textbf{Problem 27.} Let $\mathbf{u}$ and $\mathbf{v}$ be unit vectors. What is the maximum possible value of $\mathbf{u} \cdot \mathbf{v}$ and the minimum possible value? At what angles do these occur?

\textbf{Problem 28.} Given $\mathbf{a} = \langle 2, 3 \rangle$ and $\mathbf{b} = \langle k, 1 \rangle$.
\begin{enumerate}[label=(\alph*)]
    \item Find $k$ if $\mathbf{a}$ is orthogonal to $\mathbf{b}$.
    \item Find $k$ if the scalar projection of $\mathbf{b}$ onto $\mathbf{a}$ is zero.
    \item Are the answers to (a) and (b) different? Why or why not?
\end{enumerate}

\textbf{Problem 29.} Find the acute angle between the lines $2x - y = 3$ and $3x + y = 7$. (Hint: Use their normal vectors).

\textbf{Problem 30.} (Diagnostic "Find the Flaw") A student calculates the vector projection of $\mathbf{b} = \langle 1, 5 \rangle$ onto $\mathbf{a} = \langle 2, 2 \rangle$ as follows:
$$ \text{proj}_{\mathbf{a}} \mathbf{b} = \frac{\mathbf{a} \cdot \mathbf{b}}{|\mathbf{b}|^2} \mathbf{b} = \frac{12}{26} \langle 1, 5 \rangle $$
Identify the error in the formula and provide the correct calculation.

\newpage
\section*{Solutions and Comprehensive Workouts}

\textbf{1. Meaningful vs. Meaningless}
\begin{enumerate}[label=(\alph*)]
    \item \textbf{Meaningless.} $(\mathbf{u} \cdot \mathbf{v})$ is a scalar. You cannot add a scalar to a vector $\mathbf{w}$.
    \item \textbf{Meaningful (Scalar).} $|\mathbf{u}|$ is a scalar. $\mathbf{v} \cdot \mathbf{w}$ is a scalar. The product is a real number.
    \item \textbf{Meaningless.} $(\mathbf{u} \cdot \mathbf{v})$ is a scalar. You cannot compute the dot product of a scalar and a vector $\mathbf{w}$.
    \item \textbf{Meaningful (Scalar).} The sum/difference of vectors are vectors. The dot product of two vectors is a scalar.
\end{enumerate}

\textbf{2. Basic Computation}
$$ \mathbf{a} \cdot \mathbf{b} = (4)(6) + (1)(-3) + (\frac{1}{4})(-8) = 24 - 3 - 2 = 19 $$

\textbf{3. Basic Computation (i,j,k)}
$$ \mathbf{u} \cdot \mathbf{v} = (2)(1) + (-3)(2) + (5)(-1) = 2 - 6 - 5 = -9 $$

\textbf{4. Angle (2D)}
Step 1: $\mathbf{a} \cdot \mathbf{b} = (3)(5) + (4)(12) = 15 + 48 = 63$.
Step 2: $|\mathbf{a}| = \sqrt{3^2+4^2} = 5$. $|\mathbf{b}| = \sqrt{5^2+12^2} = 13$.
Step 3: $\cos \theta = \frac{63}{(5)(13)} = \frac{63}{65}$.
$$ \theta = \arccos\left(\frac{63}{65}\right) $$

\textbf{5. Angle (Parallel/Anti-parallel)}
Note that $\mathbf{b} = -3\langle 1, -2, 3 \rangle = -3\mathbf{a}$.
The vectors are parallel but point in opposite directions.
Calculation verification:
$\mathbf{a} \cdot \mathbf{b} = -3 - 12 - 27 = -42$.
$|\mathbf{a}| = \sqrt{14}$, $|\mathbf{b}| = \sqrt{9+36+81} = \sqrt{126} = 3\sqrt{14}$.
$\cos \theta = \frac{-42}{(\sqrt{14})(3\sqrt{14})} = \frac{-42}{42} = -1$.
$\theta = 180^\circ$ (or $\pi$). Significance: They are anti-parallel.

\textbf{6. Ortho/Para Check}
Check Orthogonal: $\mathbf{a} \cdot \mathbf{b} = (6)(-3) + (-2)(1) + (4)(-2) = -18 - 2 - 8 = -28 \neq 0$.
Check Parallel: Ratio check $\frac{6}{-3} = -2$, $\frac{-2}{1} = -2$, $\frac{4}{-2} = -2$.
Since ratios are constant, \textbf{Parallel}.

\textbf{7. Ortho/Para Check}
Check Orthogonal: $\mathbf{u} \cdot \mathbf{v} = (5)(6) + (3)(-8) + (7)(-2) = 30 - 24 - 14 = -8 \neq 0$.
Check Parallel: $\frac{5}{6} \neq \frac{3}{-8}$.
\textbf{Neither}.

\textbf{8. Solve for x (Orthogonal)}
Orthogonal means dot product is 0.
$\mathbf{a} \cdot \mathbf{b} = (x)(1) + (x)(x) + (2)(-3) = 0$
$x + x^2 - 6 = 0 \Rightarrow (x+3)(x-2) = 0$.
$x = -3$ or $x = 2$.

\textbf{9. Solve for b (Parallel)}
Ratios must be equal: $\frac{4}{-2} = \frac{-8}{b} = \frac{-8}{4}$.
Ratio is $-2$.
$\frac{-8}{b} = -2 \Rightarrow -2b = -8 \Rightarrow b = 4$.

\textbf{10. Triangle Angle}
Vectors originating from $A$:
$\vec{AB} = \langle 4-2, 1-(-1), 1-0 \rangle = \langle 2, 2, 1 \rangle$.
$\vec{AC} = \langle 4-2, -5-(-1), 4-0 \rangle = \langle 2, -4, 4 \rangle$.
Dot Product: $\vec{AB} \cdot \vec{AC} = (2)(2) + (2)(-4) + (1)(4) = 4 - 8 + 4 = 0$.
$\cos A = 0$, so $A = 90^\circ$.

\textbf{11. Right Triangle Check}
Sides vectors:
$\vec{PQ} = \langle 1, 3, -2 \rangle$.
$\vec{QR} = \langle 4, -2, -1 \rangle$.
$\vec{PR} = \langle 5, 1, -3 \rangle$.
Check dot products:
$\vec{PQ} \cdot \vec{QR} = 4 - 6 + 2 = 0$.
Since $\vec{PQ} \cdot \vec{QR} = 0$, the angle at $Q$ is $90^\circ$.
Yes, it is a right triangle at vertex $Q$.

\textbf{12. Direction Cosines}
$|\mathbf{v}| = \sqrt{1^2+2^2+2^2} = \sqrt{9} = 3$.
Direction Cosines:
$\cos \alpha = 1/3, \quad \cos \beta = 2/3, \quad \cos \gamma = 2/3$.
Angles:
$\alpha = \arccos(1/3) \approx 70.5^\circ$, $\beta = \arccos(2/3) \approx 48.2^\circ$, $\gamma \approx 48.2^\circ$.

\textbf{13. Rhombus Proof}
Diagonals are $\mathbf{d}_1 = \mathbf{a} + \mathbf{b}$ and $\mathbf{d}_2 = \mathbf{a} - \mathbf{b}$.
Test orthogonality:
$\mathbf{d}_1 \cdot \mathbf{d}_2 = (\mathbf{a} + \mathbf{b}) \cdot (\mathbf{a} - \mathbf{b})$.
$= \mathbf{a} \cdot \mathbf{a} - \mathbf{a} \cdot \mathbf{b} + \mathbf{b} \cdot \mathbf{a} - \mathbf{b} \cdot \mathbf{b}$.
$= |\mathbf{a}|^2 - |\mathbf{b}|^2$.
Since it is a rhombus, side lengths are equal ($|\mathbf{a}| = |\mathbf{b}|$).
$= 0$. Thus, diagonals are perpendicular.

\textbf{14. Cyclic Vectors}
Given $\mathbf{u} + \mathbf{v} = -\mathbf{w}$.
Square both sides (dot product with self):
$|\mathbf{u} + \mathbf{v}|^2 = |-\mathbf{w}|^2$
$(\mathbf{u} + \mathbf{v}) \cdot (\mathbf{u} + \mathbf{v}) = |\mathbf{w}|^2$
$|\mathbf{u}|^2 + 2(\mathbf{u} \cdot \mathbf{v}) + |\mathbf{v}|^2 = |\mathbf{w}|^2$
$3^2 + 2(\mathbf{u} \cdot \mathbf{v}) + 4^2 = 5^2$
$9 + 2(\mathbf{u} \cdot \mathbf{v}) + 16 = 25$
$25 + 2(\mathbf{u} \cdot \mathbf{v}) = 25$
$2(\mathbf{u} \cdot \mathbf{v}) = 0 \Rightarrow \mathbf{u} \cdot \mathbf{v} = 0$.

\textbf{15. Scalar Projection}
$\text{comp}_{\mathbf{a}}\mathbf{b} = \frac{\mathbf{a} \cdot \mathbf{b}}{|\mathbf{a}|}$.
$\mathbf{a} \cdot \mathbf{b} = 4 + 2 + 0 = 6$.
$|\mathbf{a}| = \sqrt{1+1+1} = \sqrt{3}$.
Result: $\frac{6}{\sqrt{3}} = 2\sqrt{3}$.

\textbf{16. Vector Projection}
$\text{proj}_{\mathbf{a}}\mathbf{b} = \left( \frac{\mathbf{a} \cdot \mathbf{b}}{|\mathbf{a}|^2} \right) \mathbf{a}$.
$\mathbf{a} \cdot \mathbf{b} = (2)(-2) + (5)(3) + (-1)(1) = -4 + 15 - 1 = 10$.
$|\mathbf{a}|^2 = 2^2 + 5^2 + (-1)^2 = 4 + 25 + 1 = 30$.
$\text{proj} = \frac{10}{30} \langle 2, 5, -1 \rangle = \frac{1}{3} \langle 2, 5, -1 \rangle = \langle \frac{2}{3}, \frac{5}{3}, -\frac{1}{3} \rangle$.

\textbf{17. Vector Decomposition}
$\mathbf{u}_{\parallel} = \text{proj}_{\mathbf{v}} \mathbf{u}$.
$\mathbf{u} \cdot \mathbf{v} = 3(1) + 4(1) = 7$. $|\mathbf{v}|^2 = 2$.
$\mathbf{u}_{\parallel} = \frac{7}{2} \langle 1, 1 \rangle = \langle 3.5, 3.5 \rangle$.
$\mathbf{u}_{\perp} = \mathbf{u} - \mathbf{u}_{\parallel} = \langle 3, 4 \rangle - \langle 3.5, 3.5 \rangle = \langle -0.5, 0.5 \rangle$.
Check: $\langle 3.5, 3.5 \rangle \cdot \langle -0.5, 0.5 \rangle = -1.75 + 1.75 = 0$. Correct.

\textbf{18. Projection Consistency}
The vector projection of $\mathbf{b}$ onto $\mathbf{a}$ must be parallel to $\mathbf{a}$.
$\mathbf{a} = \langle 4, 2 \rangle$. The result given is $\langle 2, 1 \rangle$.
Check proportionality: $\mathbf{a} = 2 \langle 2, 1 \rangle$.
Yes, the result is consistent because the projection vector is a scalar multiple of $\mathbf{a}$.

\textbf{19. Orthogonal Component}
$\text{orth}_{\mathbf{a}}\mathbf{b} = \mathbf{b} - \text{proj}_{\mathbf{a}}\mathbf{b}$.
$\mathbf{a} \cdot \mathbf{b} = 8 + 0 + 2 = 10$. $|\mathbf{a}|^2 = 4 + 0 + 1 = 5$.
$\text{proj}_{\mathbf{a}}\mathbf{b} = \frac{10}{5} \langle 2, 0, 1 \rangle = 2 \langle 2, 0, 1 \rangle = \langle 4, 0, 2 \rangle$.
$\text{orth} = \langle 4, 1, 2 \rangle - \langle 4, 0, 2 \rangle = \langle 0, 1, 0 \rangle$.

\textbf{20. Work (Vector)}
Displacement $\mathbf{d} = \vec{PQ} = \langle 3-1, 4-(-1), -1-2 \rangle = \langle 2, 5, -3 \rangle$.
$W = \mathbf{F} \cdot \mathbf{d} = (5)(2) + (2)(5) + (6)(-3) = 10 + 10 - 18 = 2$ Joules.

\textbf{21. Work (Angle)}
$W = |\mathbf{F}| |\mathbf{d}| \cos \theta$.
$W = (50)(100) \cos(60^\circ) = 5000 (0.5) = 2500$ Joules.

\textbf{22. Work Interpretation}
$\mathbf{d} = \langle 2, 2, 2 \rangle$.
$W = \mathbf{F} \cdot \mathbf{d} = (3)(2) + (-1)(2) + (1)(2) = 6 - 2 + 2 = 6$.
Positive Work.
Physical Meaning: The force is acting generally in the direction of the movement (helping the motion).

\textbf{23. Unit Vector Orthogonal to Two Vectors}
Let $\mathbf{v} = \langle x, y, z \rangle$.
1) $\mathbf{v} \cdot \mathbf{a} = x + y = 0 \Rightarrow y = -x$.
2) $\mathbf{v} \cdot \mathbf{b} = x - y + z = 0$.
Substitute (1) into (2): $x - (-x) + z = 0 \Rightarrow 2x + z = 0 \Rightarrow z = -2x$.
Let $x = 1$. Then $y = -1, z = -2$. $\mathbf{v} = \langle 1, -1, -2 \rangle$.
Magnitude $|\mathbf{v}| = \sqrt{1 + 1 + 4} = \sqrt{6}$.
Unit vector: $\frac{1}{\sqrt{6}} \langle 1, -1, -2 \rangle$.

\textbf{24. Solve for x (Angle)}
$\mathbf{a} \cdot \mathbf{b} = x$.
$|\mathbf{a}| = \sqrt{1+0+1} = \sqrt{2}$. $|\mathbf{b}| = \sqrt{x^2+1}$.
$\cos(60^\circ) = 0.5$.
$0.5 = \frac{x}{\sqrt{2}\sqrt{x^2+1}}$.
Square both sides: $0.25 = \frac{x^2}{2(x^2+1)} \Rightarrow \frac{1}{4} = \frac{x^2}{2x^2+2}$.
$2x^2 + 2 = 4x^2 \Rightarrow 2 = 2x^2 \Rightarrow x^2 = 1 \Rightarrow x = \pm 1$.
Check validity: If $x=-1$, dot product is negative, angle is obtuse ($>90$). We need $60^\circ$ (acute).
So $x = 1$. (Check: $\frac{1}{\sqrt{2}\sqrt{2}} = 1/2$. Correct).

\textbf{25. Parallelogram Proof}
LHS: $(\mathbf{a} + \mathbf{b}) \cdot (\mathbf{a} + \mathbf{b}) + (\mathbf{a} - \mathbf{b}) \cdot (\mathbf{a} - \mathbf{b})$.
$= (|\mathbf{a}|^2 + 2\mathbf{a}\cdot\mathbf{b} + |\mathbf{b}|^2) + (|\mathbf{a}|^2 - 2\mathbf{a}\cdot\mathbf{b} + |\mathbf{b}|^2)$.
$= 2|\mathbf{a}|^2 + 2|\mathbf{b}|^2$. (Q.E.D)

\textbf{26. Magnitude of Combination}
$|2\mathbf{u} - \mathbf{v}|^2 = (2\mathbf{u} - \mathbf{v}) \cdot (2\mathbf{u} - \mathbf{v})$.
$= 4|\mathbf{u}|^2 - 4(\mathbf{u} \cdot \mathbf{v}) + |\mathbf{v}|^2$.
Substitute values: $4(3^2) - 4(-10) + 5^2$.
$= 36 + 40 + 25 = 101$.
Result is $\sqrt{101}$.

\textbf{27. Unit Vector Extrema}
$\mathbf{u} \cdot \mathbf{v} = |\mathbf{u}| |\mathbf{v}| \cos \theta$.
Since they are unit vectors, $|\mathbf{u}||\mathbf{v}| = 1$.
$\mathbf{u} \cdot \mathbf{v} = \cos \theta$.
Max value: $1$ (at $\theta = 0^\circ$).
Min value: $-1$ (at $\theta = 180^\circ$).

\textbf{28. Variable k Analysis}
(a) Orthogonal: $2k + 3 = 0 \Rightarrow k = -1.5$.
(b) Scalar Proj is 0: $\frac{\mathbf{b} \cdot \mathbf{a}}{|\mathbf{a}|} = 0 \Rightarrow \mathbf{b} \cdot \mathbf{a} = 0$.
$2k + 3 = 0 \Rightarrow k = -1.5$.
(c) The answers are the same. A scalar projection of 0 means the "shadow" length is 0, implying the vectors are perpendicular.

\textbf{29. Angle Between Lines}
Normals: $\mathbf{n}_1 = \langle 2, -1 \rangle$, $\mathbf{n}_2 = \langle 3, 1 \rangle$.
$\cos \theta = \frac{|\mathbf{n}_1 \cdot \mathbf{n}_2|}{|\mathbf{n}_1||\mathbf{n}_2|}$ (Abs value for acute angle).
$\mathbf{n}_1 \cdot \mathbf{n}_2 = 6 - 1 = 5$.
$|\mathbf{n}_1| = \sqrt{5}, |\mathbf{n}_2| = \sqrt{10}$.
$\cos \theta = \frac{5}{\sqrt{5}\sqrt{10}} = \frac{5}{\sqrt{50}} = \frac{5}{5\sqrt{2}} = \frac{1}{\sqrt{2}}$.
$\theta = 45^\circ$.

\textbf{30. Find the Flaw}
\textbf{Error:} The student used the wrong vector in the projection formula. They projected onto $\mathbf{a}$ but used $|\mathbf{b}|^2$ in the denominator and multiplied by $\mathbf{b}$. The projection of $\mathbf{b}$ onto $\mathbf{a}$ must be parallel to $\mathbf{a}$.
\textbf{Correction:} Formula: $\text{proj}_{\mathbf{a}} \mathbf{b} = \frac{\mathbf{a} \cdot \mathbf{b}}{|\mathbf{a}|^2} \mathbf{a}$.
$\mathbf{a} \cdot \mathbf{b} = 12$. $|\mathbf{a}|^2 = 2^2+2^2 = 8$.
Correct Answer: $\frac{12}{8} \langle 2, 2 \rangle = \frac{3}{2} \langle 2, 2 \rangle = \langle 3, 3 \rangle$.

\newpage
\section*{Concept Check List and Question Mapping}
\begin{itemize}
    \item \textbf{1. Meaningful vs. Meaningless Operations:} Problem 1, 30.
    \item \textbf{2. Basic Dot Product Calculation:} Problems 2, 3.
    \item \textbf{3. Angle Between Vectors (Algebraic to Geometric):} Problems 4, 5, 29.
    \item \textbf{4. Orthogonality Checks/Condition:} Problems 6, 8, 13, 23, 28.
    \item \textbf{5. Parallelism Checks/Condition:} Problems 6, 7, 9.
    \item \textbf{6. Geometry (Triangles, Vertices):} Problems 10, 11, 14.
    \item \textbf{7. Direction Cosines/Angles:} Problem 12.
    \item \textbf{8. Scalar Projection:} Problem 15, 28.
    \item \textbf{9. Vector Projection:} Problem 16, 18, 30.
    \item \textbf{10. Orthogonal Decomposition/Component:} Problem 17, 19.
    \item \textbf{11. Work (Physics App):} Problems 20, 21, 22.
    \item \textbf{12. Solving for Variables (Quadratic/System):} Problems 8, 9, 23, 24.
    \item \textbf{13. Proofs/Properties:} Problems 13, 25, 26, 27.
\end{itemize}

\end{document}